\chapter[FCC Contact Microphysics]{Contact Microphysics and Homogenization of the Active FCC Vacuum Scaffold}
\label{app:fcc_microphysics}

\section{Scope, Status of Results, and Constitutive Limits}
\label{sec:fcc_microphysics_scope}

This appendix develops the baseline contact mechanics, harmonic Bloch structure, and long-wavelength homogenization of the active FCC Soft Floor. It establishes the gravity-independent FCC scaffold carrier on which the later gravitational, condensate, localized-defect, and photon-transport constitutive sectors act.

The analysis is restricted to small perturbations about a fixed active-phase equilibrium. The retained contact coefficients are tangent coefficients of that state and are not assumed to remain valid through the nonlinear FCC-to-SC phase transition. The interstitial-condensate coherent volume constraint and localized defect-induced relative-flow response are developed separately in \hyperref[app:condensate_impedance]{Appendix~\ref*{app:condensate_impedance}}.

The gravity-specific endpoint-transfer map, gravitationally referred inertia, and loaded propagation law are developed in \hyperref[app:gravitational_endpoint_transfer]{Appendix~\ref*{app:gravitational_endpoint_transfer}}, using the baseline FCC scaffold stiffness and primitive-cell inertia normalization derived here as the unstrained constitutive reference.

The mechanical sequence developed in this appendix proceeds from FCC contact geometry through objective bond kinematics, the parity-even contact action, primitive-cell assembly, the explicit harmonic Bloch stiffness, and the long-wavelength acoustic tensor.

Two nested descriptions are distinguished. The extended constitutive contact theory may contain center translation $\mathbf{u}$, mechanical microrotation $\boldsymbol{\vartheta}$, and the internal scalar phase coordinate $\varphi$. The phase-difference potential is retained because it is required in later TFT sectors. However, the present appendix does not derive an independent pristine-vacuum kinetic coefficient conjugate to $\dot{\varphi}$. A phase-difference stiffness therefore does not by itself establish an independently propagating bulk phase eigenmode.

For the acoustic analysis, the retained parent system is the six-degree-of-freedom translation--microrotation sector,
\begin{equation}
    \mathbf{q}_{\mathbf{R}}^{(6)}
    \equiv
    \begin{pmatrix}
        \mathbf{u}_{\mathbf{R}}\\
        \boldsymbol{\vartheta}_{\mathbf{R}}
    \end{pmatrix},
\end{equation}
with positive primitive-cell kinetic matrix
\begin{equation}
    \mathbf{M}_{\mathrm{cell}}^{(6)}
    =
    \begin{pmatrix}
        M_{\mathrm{D}}\mathbf{I}_3 & 0\\
        0 & I_{\mathrm{D}}\mathbf{I}_3
    \end{pmatrix},
    \qquad
    M_{\mathrm{D}}>0,
    \qquad
    I_{\mathrm{D}}>0.
\end{equation}

Within this retained acoustic sector, the explicit Bloch operator and its translational, translation--microrotation, and rotational blocks are derived below. For positive retained contact coefficients, the harmonic stiffness is positive semidefinite for every wave vector. The microrotation sector possesses a finite zone-center gap and can therefore be dynamically eliminated by a Schur complement in the acoustic limit.

The long-wavelength FCC continuum is cubic in general. Exact leading isotropy of the present nearest-neighbor translational contact law requires the constitutive specialization
\begin{equation}
    \kappa_{\mathrm{n}}=\kappa_{\mathrm{t}}.
\end{equation}
This equality is not implied by FCC geometry alone and remains a microscopic constitutive requirement of the present parent theory.

The coherent interstitial-volume constraint is not imposed in this appendix. Its discrete FCC-cell representation, pressure reaction, and continuum reduction are developed in \hyperref[app:condensate_impedance]{Appendix~\ref*{app:condensate_impedance}}. The present appendix supplies the unconstrained scaffold dynamics on which that separate constraint acts.

The status of the principal results is summarized below.

\begin{table}[h!]
\centering
\renewcommand{\arraystretch}{1.3}
\begin{tabular}{p{0.45\linewidth} p{0.5\linewidth}}
\hline
\textbf{Result} & \textbf{Status in this appendix} \\
\hline
$d_{\mathrm{nn}}=2r_{\mathrm{D}}$ & FCC geometric identity \\
$V_{\mathrm{p}}=4\sqrt{2}r_{\mathrm{D}}^{3}$ & primitive-cell geometric identity \\
baseline transverse chain dispersion & derived directional normalization branch \\
objective contact measures & derived from frame-indifferent contact kinematics \\
parity-even quadratic contact action & leading local nearest-neighbor tangent model \\
primitive-cell kinetic bookkeeping & translational and microrotational inertia assigned once per primitive cell \\
explicit six-DOF FCC Bloch operator & derived from the six undirected bonds \\
harmonic Bloch stiffness non-negativity & derived for positive retained contact coefficients \\
zone-center microrotation gap & derived \\
dynamic low-energy microrotation elimination & derived by Schur complement \\
FCC cubic elastic constants & derived in the long-wavelength limit \\
$\kappa_{\mathrm{n}}=\kappa_{\mathrm{t}}$ & condition for exact leading isotropy; constitutive specialization \\
$A_{\mathrm{ct}}^{(\mathrm{modal})}$ & constitutive surface-to-cell normalization \\
$\rho_{\mathrm{sh}}^{(\mathrm{micro},0)}$ and $K_{\mathrm{sh}}^{(\mathrm{micro},0)}$ & microscopic scaffold coefficients derived from primitive-cell normalization \\
internal phase-difference potential & retained in the extended constitutive architecture \\
independent pristine-vacuum phase inertia & open \\
selected dispersion-suppressed transverse photon branch & developed separately in \hyperref[app:photon_dispersion_relaxation]{Appendix~\ref*{app:photon_dispersion_relaxation}} \\
complete finite-$\mathbf{k}$ branch classification & open \\
\hline
\end{tabular}
\end{table}

The active Soft Floor is treated as a dual mechanical structure. The solid FCC scaffold supplies the baseline microscopic contact rigidity and primitive-cell inertia normalization. In the conservative coherent branch, the interstitial condensate is carried with the scaffold without relative flow and supplies the pressure reaction associated with the coherent volume constraint developed in \hyperref[app:condensate_impedance]{Appendix~\ref*{app:condensate_impedance}}. Entrainment fixes the relative kinematics of the two phases but does not, at the present constitutive order, establish an independently additive condensate contribution to the scaffold primitive-cell inertia normalization. The Hard Floor supplies the phase-locked structural address of the active node and is not introduced here as an independent finite-frequency oscillator.

The positive-semidefinite harmonic result established below is a tangent-state statement. It does not constitute a proof of nonlinear, prestressed, or phase-transition stability. Likewise, the acoustic Schur reduction controls the low-frequency neighborhood of $\mathbf{k}=\mathbf{0}$ and does not replace a complete classification of every finite-wave-number branch throughout the first Brillouin zone.

The present appendix therefore establishes the reference FCC scaffold contact action, objective bond kinematics, primitive-cell kinetic bookkeeping, explicit harmonic Bloch structure, dynamic microrotation reduction, and long-wavelength continuum limit. The gravitational endpoint-transfer, condensate, localized-defect, and reduced photon-transport sectors act on this baseline while preserving the objective contact kinematics and primitive-cell counting developed below.


\section{Baseline Nearest-Neighbor Transverse FCC Contact Action}
\label{sec:fcc_baseline_transverse_action}

Consider a representative nearest-neighbor sequence in the active FCC Soft Floor. The nearest-neighbor separation is
\begin{equation}
    d_{\mathrm{nn}} = 2r_{\mathrm{D}}.
\end{equation}

Let $u_n(t)$ denote the transverse displacement of the $n$-th active FCC node along one selected polarization orthogonal to the bond direction. For the contact between neighboring nodes $n$ and $n+1$, the transverse strain is
\begin{equation}
    \varepsilon_{n,n+1}^{\perp} = \frac{u_{n+1}-u_n}{d_{\mathrm{nn}}}.
\end{equation}

The baseline surface shear tension is $\tau_{\mathrm{D}}$. The corresponding transverse contact energy per unit participating interface area is
\begin{equation}
    \mathcal{U}_{n,n+1}^{\perp} = \frac{\tau_{\mathrm{D}}}{2d_{\mathrm{nn}}^2} \left( u_{n+1}-u_n \right)^2.
\end{equation}

The two-node potential Hessian is
\begin{equation}
    \mathbf{H}_U^{(n,n+1)} = \frac{\tau_{\mathrm{D}}}{d_{\mathrm{nn}}^2}
    \begin{pmatrix}
        1 & -1 \\
        -1 & 1
    \end{pmatrix}.
\end{equation}

This Hessian is positive semidefinite, with the uniform-translation mode as its null direction. The baseline scaffold surface inertia is $\mu_{\mathrm{D}}$, so the directional Lagrangian per unit participating interface area is
\begin{equation}
    \mathcal{L}_{\perp}^{(0)} = \sum_n \left[ \frac{1}{2} \mu_{\mathrm{D}} \dot{u}_n^2 - \frac{\tau_{\mathrm{D}}}{2d_{\mathrm{nn}}^2} \left( u_{n+1}-u_n \right)^2 \right].
\end{equation}

Variation with respect to $u_n$ gives
\begin{equation}
    \mu_{\mathrm{D}} \ddot{u}_n = \frac{\tau_{\mathrm{D}}}{d_{\mathrm{nn}}^2} \left( u_{n+1} - 2u_n + u_{n-1} \right).
\end{equation}

For a plane-wave mode,
\begin{equation}
    u_n(t) = u_0 \exp \left[ i \left( knd_{\mathrm{nn}} - \omega t \right) \right],
\end{equation}
the nearest-neighbor difference is
\begin{equation}
    u_{n+1} - 2u_n + u_{n-1} = -4u_n \sin^2 \left( \frac{kd_{\mathrm{nn}}}{2} \right).
\end{equation}

The baseline directional dispersion relation is therefore
\begin{equation}
    \omega^2(k) = \frac{4\tau_{\mathrm{D}}}{\mu_{\mathrm{D}}d_{\mathrm{nn}}^2} \sin^2 \left( \frac{kd_{\mathrm{nn}}}{2} \right).
\end{equation}

Using the baseline surface stiffness--inertia identity
\begin{equation}
    c_{\mathrm{E}}^2 = \frac{\tau_{\mathrm{D}}}{\mu_{\mathrm{D}}},
\end{equation}
the dispersion becomes
\begin{equation}
    \omega^2(k) = \frac{4c_{\mathrm{E}}^2}{d_{\mathrm{nn}}^2} \sin^2 \left( \frac{kd_{\mathrm{nn}}}{2} \right).
\end{equation}

In the local continuum limit,
\begin{equation}
    kd_{\mathrm{nn}} \ll 1,
\end{equation}
so
\begin{equation}
    \omega^2 = c_{\mathrm{E}}^2k^2 + \mathcal{O} \left( k^4d_{\mathrm{nn}}^2 \right).
\end{equation}

Hence the long-wavelength phase velocity is
\begin{equation}
    c_T^{(0)} = \lim_{k\rightarrow0} \frac{\omega}{k} = c_{\mathrm{E}}.
\end{equation}

On the positive-frequency monotonic branch,
\begin{equation}
    0 \leq kd_{\mathrm{nn}} \leq \pi,
\end{equation}
the dispersion and group velocity are
\begin{equation}
    \omega(k) = \frac{2c_{\mathrm{E}}}{d_{\mathrm{nn}}} \sin \left( \frac{kd_{\mathrm{nn}}}{2} \right),
\end{equation}
and
\begin{equation}
    v_g = \frac{\mathrm{d}\omega}{\mathrm{d}k} = c_{\mathrm{E}} \cos \left( \frac{kd_{\mathrm{nn}}}{2} \right),
\end{equation}
so
\begin{equation}
    \lim_{k\rightarrow0} v_g = c_{\mathrm{E}}.
\end{equation}

This one-dimensional sequence is a directional reduction used to establish the baseline transverse normalization. In the full FCC assembly, six undirected nearest-neighbor bonds contribute stiffness to the primitive-cell potential, while the generalized scaffold inertial coefficients associated with the retained node coordinates are assigned once through the primitive-cell mass matrix. The coherent condensate constraint is introduced separately in \hyperref[app:condensate_impedance]{Appendix~\ref*{app:condensate_impedance}}; condensate entrainment is not counted here as an independently additive bond or primitive-cell inertia. The directional derivation must therefore not be interpreted as assigning an independent complete nodal mass to every bond.

The nearest-neighbor branch derived here is the baseline normalization branch and is not identified with the completed physical photon eigenbranch. In particular, this scalar one-dimensional dispersion is not the full three-dimensional FCC finite-wave-number dispersion. The selected positive dispersion-suppressed transverse photon branch, its higher-gradient closure, and its observational screening are developed in \hyperref[app:photon_dispersion_relaxation]{Appendix~\ref*{app:photon_dispersion_relaxation}}.


\section{Objective Nearest-Neighbor FCC Contact Kinematics}

The directional action above captures one transverse nearest-neighbor sequence. To assemble the active Soft Floor as a three-dimensional FCC lattice, the contact law must be written in an objective form. A contact energy depending only on the relative displacement of two node centers is not sufficient, because a rigid-body rotation produces relative center motion even when the physical contact surfaces undergo no elastic deformation. The bond kinematics must therefore include nodal translations, nodal microrotations, and internal phase differences.

Consider two nearest-neighbor FCC nodes labeled $i$ and $j$. Their reference center-to-center vector is
\begin{equation}
    \mathbf{d}_{ij} = d_{\mathrm{nn}}\hat{\mathbf{n}}_{ij},
\end{equation}
where
\begin{equation}
    d_{\mathrm{nn}} = 2r_{\mathrm{D}}, \qquad |\hat{\mathbf{n}}_{ij}| = 1.
\end{equation}

The normal projector associated with the contact is
\begin{equation}
    \mathbf{P}_{\mathrm{n}}^{(ij)} = \hat{\mathbf{n}}_{ij}\hat{\mathbf{n}}_{ij}^{\mathsf{T}},
\end{equation}
and the transverse projector is
\begin{equation}
    \mathbf{P}_{\perp}^{(ij)} = \mathbf{I} - \hat{\mathbf{n}}_{ij}\hat{\mathbf{n}}_{ij}^{\mathsf{T}}.
\end{equation}

Let $\mathbf{u}_i, \mathbf{u}_j$ be the infinitesimal translations of the two node centers. Let $\boldsymbol{\vartheta}_i, \boldsymbol{\vartheta}_j$ be their infinitesimal microrotations. These are mechanical rotations of the local contact frames. They are not identical to the internal scalar phase variable. The internal phases are denoted $\varphi_i, \varphi_j$.

Thus the contact theory uses three types of local variables:
\begin{itemize}
    \item $\mathbf{u}$ for center translation,
    \item $\boldsymbol{\vartheta}$ for mechanical microrotation,
    \item $\varphi$ for internal topological phase.
\end{itemize}
They are kept distinct at the present constitutive order.

The reference contact point lies at $+r_{\mathrm{D}}\hat{\mathbf{n}}_{ij}$ relative to the center of node $i$, and at $-r_{\mathrm{D}}\hat{\mathbf{n}}_{ij}$ relative to the center of node $j$. The displacement of the contact point carried by node $i$ is
\begin{equation}
    \mathbf{u}_i^{\mathrm{c}} = \mathbf{u}_i + r_{\mathrm{D}} \boldsymbol{\vartheta}_i \times \hat{\mathbf{n}}_{ij}.
\end{equation}

The displacement of the corresponding contact point carried by node $j$ is
\begin{equation}
    \mathbf{u}_j^{\mathrm{c}} = \mathbf{u}_j - r_{\mathrm{D}} \boldsymbol{\vartheta}_j \times \hat{\mathbf{n}}_{ij}.
\end{equation}

The relative displacement of the two material contact points is therefore $\Delta\mathbf{u}_{ij}^{\mathrm{c}} = \mathbf{u}_j^{\mathrm{c}} - \mathbf{u}_i^{\mathrm{c}}$. Substitution gives
\begin{equation}
    \Delta\mathbf{u}_{ij}^{\mathrm{c}} = \mathbf{u}_j - \mathbf{u}_i - r_{\mathrm{D}} (\boldsymbol{\vartheta}_i+\boldsymbol{\vartheta}_j) \times \hat{\mathbf{n}}_{ij}.
\end{equation}
This is the objective contact-point displacement. It is the correct local deformation measure entering the nearest-neighbor contact energy.

The normal contact deformation is
\begin{equation}
    \delta_{\mathrm{n}}^{(ij)} = \hat{\mathbf{n}}_{ij} \cdot \Delta\mathbf{u}_{ij}^{\mathrm{c}}.
\end{equation}
The rotational term is transverse to $\hat{\mathbf{n}}_{ij}$, and therefore
\begin{equation}
    \delta_{\mathrm{n}}^{(ij)} = \hat{\mathbf{n}}_{ij} \cdot (\mathbf{u}_j-\mathbf{u}_i).
\end{equation}

The objective tangential slip is
\begin{equation}
    \boldsymbol{\delta}_{\mathrm{t}}^{(ij)} = \mathbf{P}_{\perp}^{(ij)} \Delta\mathbf{u}_{ij}^{\mathrm{c}}.
\end{equation}
Explicitly,
\begin{equation}
    \boldsymbol{\delta}_{\mathrm{t}}^{(ij)} = \mathbf{P}_{\perp}^{(ij)} \left[ \mathbf{u}_j-\mathbf{u}_i - r_{\mathrm{D}} (\boldsymbol{\vartheta}_i+\boldsymbol{\vartheta}_j) \times \hat{\mathbf{n}}_{ij} \right].
\end{equation}
This is the relative displacement of the physical contact surfaces in the tangent plane. It differs from the simpler expression $\mathbf{P}_{\perp}^{(ij)}(\mathbf{u}_j-\mathbf{u}_i)$, which would incorrectly assign elastic slip to a rigidly rotating pair of nodes.

Define the relative microrotation by $\Delta\boldsymbol{\vartheta}_{ij} = \boldsymbol{\vartheta}_j-\boldsymbol{\vartheta}_i$. Its transverse part gives the local rolling or bending measure,
\begin{equation}
    \boldsymbol{\delta}_{\mathrm{b}}^{(ij)} = \mathbf{P}_{\perp}^{(ij)} \Delta\boldsymbol{\vartheta}_{ij}.
\end{equation}
Its axial part gives the relative twist about the contact direction,
\begin{equation}
    \delta_{\mathrm{tw}}^{(ij)} = \hat{\mathbf{n}}_{ij} \cdot \Delta\boldsymbol{\vartheta}_{ij}.
\end{equation}

The internal phase difference across the bond is
\begin{equation}
    \Delta\varphi_{ij} = \varphi_j-\varphi_i.
\end{equation}
Only $\Delta\varphi_{ij}$, not $\varphi_i$ or $\varphi_j$ separately, may enter the contact energy. The bond action is therefore invariant under a uniform internal phase displacement,
\begin{equation}
    \varphi_i\rightarrow\varphi_i+\varphi_0, \qquad \varphi_j\rightarrow\varphi_j+\varphi_0.
\end{equation}
This is the local contact expression of the phase-shift freedom of the Soft Floor.

The objectivity of the construction can be checked directly. Consider a rigid-body motion composed of a uniform translation $\mathbf{a}$ and an infinitesimal rotation $\boldsymbol{\Omega}$. The nodal center displacements are
\begin{equation}
    \mathbf{u}_i = \mathbf{a}+\boldsymbol{\Omega}\times\mathbf{X}_i, \qquad \mathbf{u}_j = \mathbf{a}+\boldsymbol{\Omega}\times\mathbf{X}_j.
\end{equation}
Since $\mathbf{X}_j-\mathbf{X}_i = d_{\mathrm{nn}}\hat{\mathbf{n}}_{ij}$, the relative center displacement is
\begin{equation}
    \mathbf{u}_j-\mathbf{u}_i = d_{\mathrm{nn}} \boldsymbol{\Omega}\times\hat{\mathbf{n}}_{ij}.
\end{equation}

For a rigid rotation of the node frames, $\boldsymbol{\vartheta}_i = \boldsymbol{\vartheta}_j = \boldsymbol{\Omega}$. Then
\begin{equation}
    \Delta\mathbf{u}_{ij}^{\mathrm{c}} = d_{\mathrm{nn}} \boldsymbol{\Omega}\times\hat{\mathbf{n}}_{ij} - 2r_{\mathrm{D}} \boldsymbol{\Omega}\times\hat{\mathbf{n}}_{ij}.
\end{equation}
Using $d_{\mathrm{nn}}=2r_{\mathrm{D}}$, one obtains $\Delta\mathbf{u}_{ij}^{\mathrm{c}} = \mathbf{0}$.

Consequently,
\begin{equation}
    \delta_{\mathrm{n}}^{(ij)}=0, \qquad \boldsymbol{\delta}_{\mathrm{t}}^{(ij)}=\mathbf{0}.
\end{equation}
Because the two nodes undergo the same microrotation,
\begin{equation}
    \boldsymbol{\delta}_{\mathrm{b}}^{(ij)}=\mathbf{0}, \qquad \delta_{\mathrm{tw}}^{(ij)}=0.
\end{equation}
A uniform phase displacement gives $\Delta\varphi_{ij}=0$.

Thus all retained contact deformation measures vanish under rigid translation, rigid rotation, and uniform internal phase displacement. The bond law constructed from these measures is frame-indifferent to the linearized order retained here.

The resulting objective nearest-neighbor deformation set is
\begin{equation}
    \left\{ \delta_{\mathrm{n}}^{(ij)}, \boldsymbol{\delta}_{\mathrm{t}}^{(ij)}, \boldsymbol{\delta}_{\mathrm{b}}^{(ij)}, \delta_{\mathrm{tw}}^{(ij)}, \Delta\varphi_{ij} \right\}.
\end{equation}
This set is the local kinematic foundation for the parity-even quadratic bond action.


\section{Parity-Even Quadratic Bond Action}

The contact energy is now written as a minimal diagonal quadratic form built from the objective contact measures above, subject to three restrictions:
\begin{itemize}
    \item rigid-motion invariance,
    \item reciprocity of an unpolarized nearest-neighbor bond,
    \item parity evenness of the active vacuum.
\end{itemize}

Introduce the tangent contact coefficients
\begin{equation}
    \kappa_{\mathrm{n}}, \qquad \kappa_{\mathrm{t}}, \qquad \kappa_{\mathrm{b}}, \qquad \kappa_{\mathrm{tw}}, \qquad \kappa_\varphi.
\end{equation}

All coefficients are tangent coefficients of the active phase. Their dimensions are chosen so that the contact energy is expressed per unit participating interface area ($[\kappa] = \mathrm{N/m^3}$). The rotational and phase variables enter through the associated arc-length factors $r_{\mathrm{D}}\boldsymbol{\vartheta}$ and $r_{\mathrm{D}}\varphi$.

The parity-even quadratic contact energy is
\begin{equation}
\begin{aligned}
    \mathcal{U}_{ij}^{(2)} ={}& \frac{1}{2} \kappa_{\mathrm{n}} (\delta_{\mathrm{n}}^{(ij)})^2 + \frac{1}{2} \kappa_{\mathrm{t}} |\boldsymbol{\delta}_{\mathrm{t}}^{(ij)}|^2 + \frac{1}{2} \kappa_{\mathrm{b}} r_{\mathrm{D}}^2 |\boldsymbol{\delta}_{\mathrm{b}}^{(ij)}|^2 \\
    &+ \frac{1}{2} \kappa_{\mathrm{tw}} r_{\mathrm{D}}^2 (\delta_{\mathrm{tw}}^{(ij)})^2 + \frac{1}{2} \kappa_\varphi r_{\mathrm{D}}^2 (\Delta\varphi_{ij})^2.
\end{aligned}
\end{equation}

No tangential--bending cross term is included in the baseline active vacuum. The candidate scalar $\boldsymbol{\delta}_{\mathrm{t}}^{(ij)} \cdot \boldsymbol{\delta}_{\mathrm{b}}^{(ij)}$ is parity odd: $\boldsymbol{\delta}_{\mathrm{t}}$ is a polar vector, while $\boldsymbol{\delta}_{\mathrm{b}}$, being derived from microrotation, is axial. Their dot product is a pseudoscalar. It cannot appear in the parity-even constitutive energy of an unpolarized inversion-symmetric vacuum. This is why the baseline contact action strictly excludes a $\kappa_{\mathrm{tb}}$ coefficient.

A chiral extension may contain a term of the form $\chi_{ij} \kappa_{\mathrm{tb}} r_{\mathrm{D}} \boldsymbol{\delta}_{\mathrm{t}}^{(ij)} \cdot \boldsymbol{\delta}_{\mathrm{b}}^{(ij)}$, where $\chi_{ij}$ is an explicit pseudoscalar bond chirality. Such a term belongs to a polarized or chiral defect sector. It vanishes after averaging over an unpolarized active vacuum and is not part of the baseline FCC contact action.

The harmonic phase term may be understood as the tangent approximation to a periodic phase-contact potential,
\begin{equation}
    \mathcal{U}_{\varphi}^{(ij)} = \kappa_\varphi r_{\mathrm{D}}^2 \left[ 1-\cos(\Delta\varphi_{ij}) \right].
\end{equation}
For $|\Delta\varphi_{ij}| \ll 1$, this becomes
\begin{equation}
    \mathcal{U}_{\varphi}^{(ij)} = \frac{1}{2} \kappa_\varphi r_{\mathrm{D}}^2 (\Delta\varphi_{ij})^2 + \mathcal{O}(\Delta\varphi_{ij}^4).
\end{equation}

The local stability conditions of the uncoupled parity-even bond are
\begin{equation}
    \kappa_{\mathrm{n}}>0, \qquad \kappa_{\mathrm{t}}>0, \qquad \kappa_{\mathrm{b}}>0, \qquad \kappa_{\mathrm{tw}}>0, \qquad \kappa_\varphi>0.
\end{equation}

It is useful to define the contact-strain vector
\begin{equation}
    \boldsymbol{\xi}_{ij} =
    \begin{pmatrix}
        \delta_{\mathrm{n}}^{(ij)} \\
        \boldsymbol{\delta}_{\mathrm{t}}^{(ij)} \\
        r_{\mathrm{D}}\boldsymbol{\delta}_{\mathrm{b}}^{(ij)} \\
        r_{\mathrm{D}}\delta_{\mathrm{tw}}^{(ij)} \\
        r_{\mathrm{D}}\Delta\varphi_{ij}
    \end{pmatrix}.
\end{equation}

With the tangential and bending entries understood as two-component vectors in the contact tangent plane, the bond energy can be written compactly as $\mathcal{U}_{ij}^{(2)} = \frac{1}{2} \boldsymbol{\xi}_{ij}^{\mathsf{T}} \mathbf{H}_{\mathrm{c}} \boldsymbol{\xi}_{ij}$. The baseline parity-even contact tangent matrix is block diagonal:
\begin{equation}
    \mathbf{H}_{\mathrm{c}} =
    \begin{pmatrix}
        \kappa_{\mathrm{n}} & 0 & 0 & 0 & 0 \\
        0 & \kappa_{\mathrm{t}}\mathbf{I}_2 & 0 & 0 & 0 \\
        0 & 0 & \kappa_{\mathrm{b}}\mathbf{I}_2 & 0 & 0 \\
        0 & 0 & 0 & \kappa_{\mathrm{tw}} & 0 \\
        0 & 0 & 0 & 0 & \kappa_\varphi
    \end{pmatrix}.
\end{equation}

Let the generalized two-node coordinate be
\begin{equation}
    \mathbf{q}_{ij} =
    \begin{pmatrix}
        \mathbf{u}_i \\
        \boldsymbol{\vartheta}_i \\
        \varphi_i \\
        \mathbf{u}_j \\
        \boldsymbol{\vartheta}_j \\
        \varphi_j
    \end{pmatrix}.
\end{equation}

Every component of $\boldsymbol{\xi}_{ij}$ is linear in the generalized coordinates, so $\boldsymbol{\xi}_{ij} = \mathbf{B}_{ij}\mathbf{q}_{ij}$, where $\mathbf{B}_{ij}$ is the objective bond-kinematic operator. The contact energy per unit participating area becomes
\begin{equation}
    \mathcal{U}_{ij}^{(2)} = \frac{1}{2} \mathbf{q}_{ij}^{\mathsf{T}} \overline{\mathbf{K}}_{ij} \mathbf{q}_{ij},
\end{equation}
with
\begin{equation}
    \overline{\mathbf{K}}_{ij} = \mathbf{B}_{ij}^{\mathsf{T}} \mathbf{H}_{\mathrm{c}} \mathbf{B}_{ij}.
\end{equation}
The overbar denotes a generalized stiffness block per unit participating interface area. The integrated bond stiffness entering the primitive-cell assembly is obtained only after multiplication by the effective modal area.

The bond is reciprocal under endpoint interchange ($i \leftrightarrow j$). The squared magnitudes in the energy are unchanged, and because the phase energy is even in $\Delta\varphi_{ij}$, the complete bond energy satisfies $\mathcal{U}_{ij}^{(2)} = \mathcal{U}_{ji}^{(2)}$.

For a mode with no normal deformation, no relative microrotation, and no phase difference, the contact energy reduces to
\begin{equation}
    \mathcal{U}_{ij}^{(2)} = \frac{1}{2} \kappa_{\mathrm{t}} |\boldsymbol{\delta}_{\mathrm{t}}^{(ij)}|^2.
\end{equation}
For the restricted directional transverse mode, consistency requires
\begin{equation}
    \kappa_{\mathrm{t}} = \frac{\tau_{\mathrm{D}}}{d_{\mathrm{nn}}^2}.
\end{equation}
Under the isotropic contact specialization,
\begin{equation}
    \kappa_{\mathrm{n}}=\kappa_{\mathrm{t}}\equiv\kappa_{\mathrm{c}},
\end{equation}
the common translational coefficient becomes
\begin{equation}
    \kappa_{\mathrm{c}} = \frac{\tau_{\mathrm{D}}}{d_{\mathrm{nn}}^2}.
\end{equation}
This recovers the scalar transverse contact law as one polarization of the objective three-dimensional nearest-neighbor action.

The internal phase-difference potential belongs to the extended TFT constitutive architecture. However, its presence in the potential does not by itself define a propagating bulk phase mode. Such a mode additionally requires an independently derived positive kinetic coefficient conjugate to $\dot{\varphi}$, which is not supplied by the present appendix.

For the translation--microrotation acoustic problem, define the reduced contact-strain vector
\begin{equation}
    \boldsymbol{\xi}_{ij}^{(6)} =
    \begin{pmatrix}
        \delta_{\mathrm{n}}^{(ij)} \\
        \boldsymbol{\delta}_{\mathrm{t}}^{(ij)} \\
        r_{\mathrm{D}}\boldsymbol{\delta}_{\mathrm{b}}^{(ij)} \\
        r_{\mathrm{D}}\delta_{\mathrm{tw}}^{(ij)}
    \end{pmatrix}.
\end{equation}

For the Bloch calculation it is convenient to represent the tangential and bending vectors as projected three-component ambient vectors. The corresponding reduced contact Hessian is
\begin{equation}
    \mathbf{H}_{\mathrm{c}}^{(6)}
    =
    \operatorname{diag}
    \left(
        \kappa_{\mathrm{n}},
        \kappa_{\mathrm{t}}\mathbf{I}_3,
        \kappa_{\mathrm{b}}\mathbf{I}_3,
        \kappa_{\mathrm{tw}}
    \right).
\end{equation}
The projectors contained in the bond kinematics restrict the nominal three-component tangential and bending entries to their physical two-dimensional tangent subspaces; this ambient representation therefore introduces no additional physical degree of freedom.


\section[Primitive-Cell Inertia]{Primitive-Cell Kinetic Bookkeeping}
\label{sec:fcc_cell_inertia}

The bond law above is written per unit participating interface area. Let $A_{\mathrm{ct}}$ denote the effective modal area used to integrate that surface law. For the $a$-th undirected nearest-neighbor contact,
\begin{equation}
    U_a = A_{\mathrm{ct}}\mathcal{U}_a^{(2)},
\end{equation}
and the corresponding integrated generalized stiffness block is
\begin{equation}
    \mathbf{K}_a = A_{\mathrm{ct}}\overline{\mathbf{K}}_a.
\end{equation}
The modal interpretation and normalization of $A_{\mathrm{ct}}$ are established later in this appendix.

The six undirected nearest-neighbor contacts contribute potential energy once per primitive cell. The kinetic inertia, however, must not be assigned independently to each contact as a separate copy of the primitive-node inertia. Translational and microrotational inertia belong to the retained primitive-cell coordinates and are assigned once.

For the minimal acoustic parent system,
\begin{equation}
    \mathbf{q}_{\mathbf{R}}^{(6)}
    =
    \begin{pmatrix}
        \mathbf{u}_{\mathbf{R}}\\
        \boldsymbol{\vartheta}_{\mathbf{R}}
    \end{pmatrix},
\end{equation}
the lattice bookkeeping is
\begin{equation}
    U^{(6)} = \sum_{\mathbf{R}} \sum_{a=1}^{6} U_a^{(6)}\!\left( \mathbf{q}_{\mathbf{R}}^{(6)}, \mathbf{q}_{\mathbf{R}+\mathbf{d}_a}^{(6)} \right),
\end{equation}
whereas
\begin{equation}
    T^{(6)} = \frac{1}{2} \sum_{\mathbf{R}} \dot{\mathbf{q}}_{\mathbf{R}}^{(6)\,\mathsf{T}} \mathbf{M}_{\mathrm{cell}}^{(6)} \dot{\mathbf{q}}_{\mathbf{R}}^{(6)}.
\end{equation}

The retained primitive-cell kinetic matrix is
\begin{equation}
    \mathbf{M}_{\mathrm{cell}}^{(6)}
    =
    \begin{pmatrix}
        M_{\mathrm{D}}\mathbf{I}_3 & 0\\
        0 & I_{\mathrm{D}}\mathbf{I}_3
    \end{pmatrix},
\end{equation}
with
\begin{equation}
    M_{\mathrm{D}}>0, \qquad I_{\mathrm{D}}>0.
\end{equation}

The translational inertia $M_{\mathrm{D}}$ and microrotational inertia $I_{\mathrm{D}}$ are therefore counted once per primitive cell. The microrotational inertia is left symbolic in the present appendix; its numerical microscopic value is not required for the leading acoustic reduction, provided it is finite and positive.

The broader TFT constitutive coordinate may additionally contain the internal phase,
\begin{equation}
    \mathbf{q}_{\mathbf{R}}^{(\mathrm{ext})}
    =
    \begin{pmatrix}
        \mathbf{u}_{\mathbf{R}}\\
        \boldsymbol{\vartheta}_{\mathbf{R}}\\
        \varphi_{\mathbf{R}}
    \end{pmatrix}.
\end{equation}
The phase coordinate remains part of the potential architecture because phase differences are used in later TFT sectors. However, no independent pristine-vacuum phase inertia is established here. Consequently, no full seven-component generalized eigenvalue determinant is asserted at this stage. A future microscopic derivation of a positive phase inertia would enlarge the kinetic matrix and introduce an additional dynamical sector that must then be classified independently.

No independently additive condensate inertia is included in this primitive-cell scaffold normalization. The coherent condensate kinematics and volume constraint are introduced separately in \hyperref[app:condensate_impedance]{Appendix~\ref*{app:condensate_impedance}}.

After Bloch reduction at fixed $\mathbf{k}$, the retained acoustic quadratic action has the form
\begin{equation}
    S_{\mathbf{k}}^{(6)} = \int \left[ \frac{1}{2} \dot{\mathbf{q}}_0^{\dagger} \mathbf{M}_{\mathrm{cell}}^{(6)} \dot{\mathbf{q}}_0 - \frac{1}{2} \mathbf{q}_0^{\dagger} \mathbf{K}_{\mathrm{FCC}}^{(6)}(\mathbf{k}) \mathbf{q}_0 \right] \mathrm{d}t.
\end{equation}

This displays the required assembly logic: the six bonds contribute integrated stiffness, whereas the primitive cell carries the retained generalized inertia once. Any later coupled scaffold--condensate kinetic extension must act through a consistent enlargement or coordinate transformation of the primitive-cell kinetic structure rather than by duplicating the nodal inertia bond by bond.


\section{Nearest-Neighbor FCC Bloch Assembly and Harmonic Spectrum}
\label{sec:fcc_bloch_assembly}

The active Soft Floor is an FCC Bravais lattice with one generalized node per primitive cell. Each node has twelve nearest neighbors. Because opposite directions represent the same undirected contact, the primitive-cell potential energy is assembled from six independently counted nearest-neighbor bonds.

The nearest-neighbor separation is
\begin{equation}
    d_{\mathrm{nn}}=2r_{\mathrm{D}}.
\end{equation}
A convenient set of six undirected nearest-neighbor vectors is
\begin{equation}
\begin{aligned}
    \mathbf{d}_1&=\frac{d_{\mathrm{nn}}}{\sqrt{2}}(1,1,0), & \mathbf{d}_2&=\frac{d_{\mathrm{nn}}}{\sqrt{2}}(1,-1,0),\\
    \mathbf{d}_3&=\frac{d_{\mathrm{nn}}}{\sqrt{2}}(1,0,1), & \mathbf{d}_4&=\frac{d_{\mathrm{nn}}}{\sqrt{2}}(1,0,-1),\\
    \mathbf{d}_5&=\frac{d_{\mathrm{nn}}}{\sqrt{2}}(0,1,1), & \mathbf{d}_6&=\frac{d_{\mathrm{nn}}}{\sqrt{2}}(0,1,-1).
\end{aligned}
\end{equation}
The corresponding unit vectors are $\hat{\mathbf{n}}_a=\mathbf{d}_a/d_{\mathrm{nn}}$. The negatives of these six directions supply the remaining six directed neighbors.

The six undirected directions satisfy
\begin{equation}
    \sum_{a=1}^{6} n_i^{(a)}n_j^{(a)} = 2\delta_{ij}.
\end{equation}

The conventional cubic FCC cell has edge length
\begin{equation}
    a_{\mathrm{FCC}}=2\sqrt{2}\,r_{\mathrm{D}},
\end{equation}
and the primitive-cell volume is
\begin{equation}
    V_{\mathrm{p}}=\frac{a_{\mathrm{FCC}}^3}{4}=4\sqrt{2}\,r_{\mathrm{D}}^3.
\end{equation}

For a Bloch mode of the retained translation--microrotation parent system,
\begin{equation}
    \mathbf{q}_{\mathbf{R}}^{(6)}(t)=\mathbf{q}_0 \exp[i(\mathbf{k}\cdot\mathbf{R}-\omega t)],
\end{equation}
so
\begin{equation}
    \mathbf{q}_{\mathbf{R}+\mathbf{d}_a}^{(6)} = e^{i\mathbf{k}\cdot\mathbf{d}_a} \mathbf{q}_{\mathbf{R}}^{(6)}.
\end{equation}

Define
\begin{equation}
    q_a = \mathbf{k}\cdot\mathbf{d}_a, \qquad z_a=e^{iq_a}, \qquad \alpha_a=z_a-1, \qquad \gamma_a=z_a+1.
\end{equation}
Also define
\begin{equation}
    \mathbf{P}_a = \mathbf{I}-\hat{\mathbf{n}}_a\hat{\mathbf{n}}_a^{\mathsf{T}},
\end{equation}
and the cross-product operator
\begin{equation}
    \mathbf{N}_a\mathbf{x} = \hat{\mathbf{n}}_a\times\mathbf{x}.
\end{equation}

The reduced objective Bloch contact operator is
\begin{equation}
    \boxed{
    \mathbf{B}_a^{(6)}(\mathbf{k})
    =
    \begin{pmatrix}
        \alpha_a\hat{\mathbf{n}}_a^{\mathsf{T}} & 0\\
        \alpha_a\mathbf{P}_a & r_{\mathrm{D}}\gamma_a\mathbf{N}_a\\
        0 & r_{\mathrm{D}}\alpha_a\mathbf{P}_a\\
        0 & r_{\mathrm{D}}\alpha_a\hat{\mathbf{n}}_a^{\mathsf{T}}
    \end{pmatrix}.
    }
\end{equation}

For contact $a$,
\begin{equation}
    \boldsymbol{\xi}_a^{(6)} = \mathbf{B}_a^{(6)}(\mathbf{k})\mathbf{q}_0.
\end{equation}

With the reduced contact tangent matrix $\mathbf{H}_{\mathrm{c}}^{(6)}$, the primitive-cell Bloch stiffness is
\begin{equation}
    \boxed{
    \mathbf{K}_{\mathrm{FCC}}^{(6)}(\mathbf{k})
    =
    A_{\mathrm{ct}}
    \sum_{a=1}^{6}
    \mathbf{B}_a^{(6)\dagger}(\mathbf{k})
    \mathbf{H}_{\mathrm{c}}^{(6)}
    \mathbf{B}_a^{(6)}(\mathbf{k}).
    }
\end{equation}

For arbitrary generalized amplitude $\mathbf{q}_0$,
\begin{equation}
\begin{aligned}
    \mathbf{q}_0^{\dagger}\mathbf{K}_{\mathrm{FCC}}^{(6)}\mathbf{q}_0
    =
    A_{\mathrm{ct}}\sum_{a=1}^{6}
    \left(\mathbf{B}_a^{(6)}\mathbf{q}_0\right)^{\dagger}
    \mathbf{H}_{\mathrm{c}}^{(6)}
    \left(\mathbf{B}_a^{(6)}\mathbf{q}_0\right)
    \geq 0.
\end{aligned}
\end{equation}

Hence,
\begin{equation}
    \boxed{
    \mathbf{K}_{\mathrm{FCC}}^{(6)}(\mathbf{k})\succeq0
    \qquad
    \text{for all }\mathbf{k},
    }
\end{equation}
provided
\begin{equation}
    \kappa_{\mathrm{n}}>0, \qquad
    \kappa_{\mathrm{t}}>0, \qquad
    \kappa_{\mathrm{b}}>0, \qquad
    \kappa_{\mathrm{tw}}>0.
\end{equation}

Because $\mathbf{M}_{\mathrm{cell}}^{(6)}>0$, the retained generalized eigenvalue problem
\begin{equation}
    \det\!\left[
        \mathbf{K}_{\mathrm{FCC}}^{(6)}(\mathbf{k})
        -
        \omega^2\mathbf{M}_{\mathrm{cell}}^{(6)}
    \right]
    =
    0
\end{equation}
satisfies
\begin{equation}
    \omega_s^2(\mathbf{k})\geq0
\end{equation}
for every branch $s$ of the retained six-degree-of-freedom harmonic system.

This proves harmonic spectral non-negativity of the retained reference lattice. It does not constitute a nonlinear or preloaded-state stability proof.

The translational block is
\begin{equation}
    \mathbf{K}_{uu}
    =
    4A_{\mathrm{ct}}
    \sum_{a=1}^{6}
    \sin^2\!\left(\frac{q_a}{2}\right)
    \left[
        \kappa_{\mathrm{t}}\mathbf{I}
        +
        \left(
            \kappa_{\mathrm{n}}-\kappa_{\mathrm{t}}
        \right)
        \hat{\mathbf{n}}_a\hat{\mathbf{n}}_a^{\mathsf{T}}
    \right].
\end{equation}

The translation--microrotation block is
\begin{equation}
    \mathbf{K}_{u\vartheta}
    =
    -2iA_{\mathrm{ct}}\kappa_{\mathrm{t}}r_{\mathrm{D}}
    \sum_{a=1}^{6}
    \sin q_a\,\mathbf{N}_a,
    \qquad
    \mathbf{K}_{\vartheta u}
    =
    \mathbf{K}_{u\vartheta}^{\dagger}.
\end{equation}

The rotational block is
\begin{equation}
\begin{aligned}
    \mathbf{K}_{\vartheta\vartheta}
    =
    4A_{\mathrm{ct}}r_{\mathrm{D}}^2
    \sum_{a=1}^{6}
    \Bigg[
        \kappa_{\mathrm{t}}
        \cos^2\!\left(\frac{q_a}{2}\right)\mathbf{P}_a
        +
        \sin^2\!\left(\frac{q_a}{2}\right)
        \left(
            \kappa_{\mathrm{b}}\mathbf{P}_a
            +
            \kappa_{\mathrm{tw}}
            \hat{\mathbf{n}}_a\hat{\mathbf{n}}_a^{\mathsf{T}}
        \right)
    \Bigg].
\end{aligned}
\end{equation}

At $\mathbf{k}=\mathbf{0}$,
\begin{equation}
    \mathbf{K}_{uu}(0)=0,
    \qquad
    \mathbf{K}_{u\vartheta}(0)=0,
\end{equation}
while
\begin{equation}
    \mathbf{K}_{\vartheta\vartheta}(0)
    =
    4A_{\mathrm{ct}}\kappa_{\mathrm{t}}r_{\mathrm{D}}^2
    \sum_{a=1}^{6}\mathbf{P}_a.
\end{equation}

Using
\begin{equation}
    \sum_{a=1}^{6}
    \hat{\mathbf{n}}_a\hat{\mathbf{n}}_a^{\mathsf{T}}
    =
    2\mathbf{I},
\end{equation}
one has
\begin{equation}
    \sum_{a=1}^{6}\mathbf{P}_a
    =
    4\mathbf{I},
\end{equation}
and hence
\begin{equation}
    \boxed{
    \mathbf{K}_{\vartheta\vartheta}(0)
    =
    16A_{\mathrm{ct}}\kappa_{\mathrm{t}}r_{\mathrm{D}}^2\mathbf{I}.
    }
\end{equation}

The microrotational triplet therefore has the finite zone-center frequency
\begin{equation}
    \boxed{
    \omega_{\vartheta}^2(0)
    =
    \frac{
        16A_{\mathrm{ct}}\kappa_{\mathrm{t}}r_{\mathrm{D}}^2
    }{
        I_{\mathrm{D}}
    }.
    }
\end{equation}

Thus the rotational sector is optical rather than acoustic.

A uniform rigid-body rotation remains a zero-deformation motion of a finite body or affine continuum field, but it is not represented by a spatially uniform periodic Bloch amplitude because the compatible center displacement $\mathbf{u}=\boldsymbol{\Omega}\times\mathbf{x}$ varies linearly with position. It therefore appears in the affine $\mathbf{k}\rightarrow\mathbf{0}$ objectivity limit rather than as an additional uniform Bloch zero mode.

For completeness, the decoupled harmonic potential block of the extended internal phase sector is
\begin{equation}
    K_{\varphi\varphi}(\mathbf{k})
    =
    4A_{\mathrm{ct}}\kappa_{\varphi}r_{\mathrm{D}}^2
    \sum_{a=1}^{6}
    \sin^2\!\left(\frac{q_a}{2}\right).
\end{equation}

Therefore,
\begin{equation}
    K_{\varphi\varphi}(0)=0,
\end{equation}
as required by uniform phase-shift invariance. This potential null direction does not establish a propagating phase acoustic mode because no independent phase kinetic coefficient is supplied by the present appendix. A future positive phase inertia would enlarge the dynamical problem and must then be analyzed explicitly.

The positive centered-second-difference contribution with coefficient $1/12$ used in the reduced photon action of \hyperref[app:photon_dispersion_relaxation]{Appendix~\ref*{app:photon_dispersion_relaxation}} is a higher-gradient constitutive augmentation of the selected transverse branch. It is not derived by the nearest-neighbor parent action above alone.


\section{FCC Long-Wavelength Homogenization and Cubic Elasticity}

The leading continuum tensor is obtained from a spatially affine displacement field $\mathbf{u}(\mathbf{x}) = \mathbf{H}\mathbf{x}$, where $\mathbf{H}=\nabla\mathbf{u}$. Decompose $\mathbf{H}$ into symmetric and antisymmetric parts:
\begin{equation}
    \mathbf{H} = \boldsymbol{\varepsilon}+\mathbf{W},
\end{equation}
with $\boldsymbol{\varepsilon} = \frac{1}{2}(\mathbf{H}+\mathbf{H}^{\mathsf{T}})$ and $\mathbf{W} = \frac{1}{2}(\mathbf{H}-\mathbf{H}^{\mathsf{T}})$. The symmetric tensor $\boldsymbol{\varepsilon}$ is the infinitesimal strain. The antisymmetric tensor $\mathbf{W}$ represents local rigid rotation. Let $\boldsymbol{\omega}$ be the axial vector associated with $\mathbf{W}$, so that $\mathbf{W}\hat{\mathbf{n}} = \boldsymbol{\omega}\times\hat{\mathbf{n}}$.

Across contact $a$, the affine relative center displacement is $\mathbf{u}(\mathbf{x}+\mathbf{d}_a)-\mathbf{u}(\mathbf{x}) = d_{\mathrm{nn}} \mathbf{H}\hat{\mathbf{n}}_a$. For a spatially uniform microrotation $\boldsymbol{\vartheta}$, the objective relative contact-point displacement becomes
\begin{equation}
    \Delta\mathbf{u}_a^{\mathrm{c}} = d_{\mathrm{nn}} \left[ \boldsymbol{\varepsilon}\hat{\mathbf{n}}_a + (\boldsymbol{\omega}-\boldsymbol{\vartheta}) \times \hat{\mathbf{n}}_a \right].
\end{equation}

The normal deformation is therefore $\delta_{\mathrm{n}}^{(a)} = d_{\mathrm{nn}} \hat{\mathbf{n}}_a \cdot \boldsymbol{\varepsilon} \hat{\mathbf{n}}_a$. The tangential deformation is
\begin{equation}
    \boldsymbol{\delta}_{\mathrm{t}}^{(a)} = d_{\mathrm{nn}} \mathbf{P}_\perp^{(a)} \left[ \boldsymbol{\varepsilon}\hat{\mathbf{n}}_a + (\boldsymbol{\omega}-\boldsymbol{\vartheta}) \times \hat{\mathbf{n}}_a \right],
\end{equation}
where $\mathbf{P}_\perp^{(a)} = \mathbf{I}-\hat{\mathbf{n}}_a\hat{\mathbf{n}}_a^{\mathsf{T}}$.

The six FCC directions satisfy the second-order identity $\sum_{a=1}^{6} n_i^{(a)}n_j^{(a)} = 2\delta_{ij}$. For a symmetric strain tensor, this identity implies that the summed cross term between $\boldsymbol{\varepsilon}\hat{\mathbf{n}}_a$ and $(\boldsymbol{\omega}-\boldsymbol{\vartheta})\times\hat{\mathbf{n}}_a$ vanishes. The relative-rotation contribution is positive and proportional to $|\boldsymbol{\omega}-\boldsymbol{\vartheta}|^2$.

Minimizing the affine energy with respect to the uniform microrotation gives $\boldsymbol{\vartheta}=\boldsymbol{\omega}$. Thus the microrotation follows the macroscopic rigid rotation. After relaxation,
\begin{equation}
    \boldsymbol{\delta}_{\mathrm{t}}^{(a)} = d_{\mathrm{nn}} \mathbf{P}_\perp^{(a)} \boldsymbol{\varepsilon}\hat{\mathbf{n}}_a.
\end{equation}

For a homogeneous affine deformation, neighboring cells share the same microrotation and internal phase. Therefore, $\boldsymbol{\delta}_{\mathrm{b}}^{(a)}=\mathbf{0}$, $\delta_{\mathrm{tw}}^{(a)}=0$, and $\Delta\varphi_a=0$. The bending, twisting, and phase sectors do not contribute to the leading affine elastic tensor.

Let $A_{\mathrm{ct}}$ denote the effective modal area measure multiplying the contact energy per unit interface area. The leading scaffold strain-energy density is
\begin{equation}
    W_{\mathrm{FCC}} = \frac{A_{\mathrm{ct}}}{V_{\mathrm{p}}} \sum_{a=1}^{6} \left[ \frac{1}{2}\kappa_{\mathrm{n}} (\delta_{\mathrm{n}}^{(a)})^2 + \frac{1}{2}\kappa_{\mathrm{t}} |\boldsymbol{\delta}_{\mathrm{t}}^{(a)}|^2 \right].
\end{equation}

Using the relaxed affine contact kinematics,
\begin{equation}
    |\boldsymbol{\delta}_{\mathrm{t}}^{(a)}|^2 = d_{\mathrm{nn}}^2 \left[ |\boldsymbol{\varepsilon}\hat{\mathbf{n}}_a|^2 - (\hat{\mathbf{n}}_a\cdot \boldsymbol{\varepsilon} \hat{\mathbf{n}}_a)^2 \right].
\end{equation}

Therefore,
\begin{equation}
    W_{\mathrm{FCC}} = \frac{A_{\mathrm{ct}}d_{\mathrm{nn}}^2}{2V_{\mathrm{p}}} \sum_{a=1}^{6} \left[ \kappa_{\mathrm{t}} |\boldsymbol{\varepsilon}\hat{\mathbf{n}}_a|^2 + (\kappa_{\mathrm{n}}-\kappa_{\mathrm{t}}) (\hat{\mathbf{n}}_a\cdot \boldsymbol{\varepsilon} \hat{\mathbf{n}}_a)^2 \right].
\end{equation}

The required FCC directional sums are $\sum_{a=1}^{6} |\boldsymbol{\varepsilon}\hat{\mathbf{n}}_a|^2 = 2\boldsymbol{\varepsilon}:\boldsymbol{\varepsilon}$, and
\begin{equation}
\begin{aligned}
    \sum_{a=1}^{6} (\hat{\mathbf{n}}_a\cdot \boldsymbol{\varepsilon} \hat{\mathbf{n}}_a)^2 ={}& \varepsilon_{xx}^2 + \varepsilon_{yy}^2 + \varepsilon_{zz}^2 \\
    & + \varepsilon_{xx}\varepsilon_{yy} + \varepsilon_{xx}\varepsilon_{zz} + \varepsilon_{yy}\varepsilon_{zz} \\
    & + 2(\varepsilon_{xy}^2+\varepsilon_{xz}^2+\varepsilon_{yz}^2).
\end{aligned}
\end{equation}

Substitution gives
\begin{equation}
\begin{aligned}
    W_{\mathrm{FCC}} = \frac{A_{\mathrm{ct}}d_{\mathrm{nn}}^2}{2V_{\mathrm{p}}} \Big[ & (\kappa_{\mathrm{n}}+\kappa_{\mathrm{t}}) (\varepsilon_{xx}^2+\varepsilon_{yy}^2+\varepsilon_{zz}^2) \\
    & + (\kappa_{\mathrm{n}}-\kappa_{\mathrm{t}}) (\varepsilon_{xx}\varepsilon_{yy} +\varepsilon_{xx}\varepsilon_{zz} +\varepsilon_{yy}\varepsilon_{zz}) \\
    & + 2(\kappa_{\mathrm{n}}+\kappa_{\mathrm{t}}) (\varepsilon_{xy}^2+\varepsilon_{xz}^2+\varepsilon_{yz}^2) \Big].
\end{aligned}
\end{equation}

The strain-energy density of a cubic continuum is
\begin{equation}
\begin{aligned}
    W_{\mathrm{cub}} ={}& \frac{1}{2} C_{11} (\varepsilon_{xx}^2+\varepsilon_{yy}^2+\varepsilon_{zz}^2) \\
    & + C_{12} (\varepsilon_{xx}\varepsilon_{yy} +\varepsilon_{xx}\varepsilon_{zz} +\varepsilon_{yy}\varepsilon_{zz}) \\
    & + 2C_{44} (\varepsilon_{xy}^2+\varepsilon_{xz}^2+\varepsilon_{yz}^2).
\end{aligned}
\end{equation}

Matching terms gives the leading FCC scaffold constants:
\begin{equation}
    C_{11}^{(\mathrm{FCC})} = \frac{A_{\mathrm{ct}}d_{\mathrm{nn}}^2}{V_{\mathrm{p}}} (\kappa_{\mathrm{n}}+\kappa_{\mathrm{t}}),
\end{equation}
\begin{equation}
    C_{12}^{(\mathrm{FCC})} = \frac{A_{\mathrm{ct}}d_{\mathrm{nn}}^2}{2V_{\mathrm{p}}} (\kappa_{\mathrm{n}}-\kappa_{\mathrm{t}}),
\end{equation}
\begin{equation}
    C_{44}^{(\mathrm{FCC})} = \frac{A_{\mathrm{ct}}d_{\mathrm{nn}}^2}{2V_{\mathrm{p}}} (\kappa_{\mathrm{n}}+\kappa_{\mathrm{t}}).
\end{equation}

These constants define a cubic continuum. FCC geometry alone guarantees the cubic tensor structure, not exact isotropy. Exact isotropy requires a further relation among the tangent contact coefficients.

The leading cubic continuum is mechanically stable when
\begin{equation}
    C_{44}>0, \qquad C_{11}-C_{12}>0, \qquad C_{11}+2C_{12}>0.
\end{equation}
For the present FCC coefficients, these conditions are equivalent to
\begin{equation}
    \kappa_{\mathrm{n}}>0, \qquad
    \kappa_{\mathrm{n}}+\kappa_{\mathrm{t}}>0, \qquad
    \kappa_{\mathrm{n}}+3\kappa_{\mathrm{t}}>0.
\end{equation}
The stronger bond-level positivity conditions
\begin{equation}
    \kappa_{\mathrm{n}}>0, \qquad \kappa_{\mathrm{t}}>0
\end{equation}
are sufficient for these affine continuum inequalities and are already required by the positive local bond energy.


\subsection{Dynamic Low-Energy Elimination of Microrotation}
\label{sec:fcc_dynamic_microrotation}

The affine minimization above establishes the relaxed long-wavelength elastic tensor. A propagating wave, however, is dynamical. The same relaxed tensor must therefore emerge from the coupled translation--microrotation eigenproblem in the acoustic limit.

Under the isotropic contact specialization
\begin{equation}
    \kappa_{\mathrm{n}}=\kappa_{\mathrm{t}}\equiv\kappa_{\mathrm{c}},
\end{equation}
the leading nonvanishing small-$\mathbf{k}$ blocks of the FCC stiffness are
\begin{equation}
    \mathbf{K}_{uu}^{(2)}
    =
    8A_{\mathrm{ct}}\kappa_{\mathrm{c}}r_{\mathrm{D}}^2k^2\mathbf{I},
\end{equation}
\begin{equation}
    \mathbf{K}_{u\vartheta}^{(1)}
    =
    -8iA_{\mathrm{ct}}\kappa_{\mathrm{c}}r_{\mathrm{D}}^2[\mathbf{k}]_{\times},
\end{equation}
and
\begin{equation}
    \mathbf{K}_{\vartheta\vartheta}^{(0)}
    =
    16A_{\mathrm{ct}}\kappa_{\mathrm{c}}r_{\mathrm{D}}^2\mathbf{I}.
\end{equation}

Here $[\mathbf{k}]_{\times}$ is the antisymmetric operator satisfying
\begin{equation}
    [\mathbf{k}]_{\times}\mathbf{x}
    =
    \mathbf{k}\times\mathbf{x}.
\end{equation}

The rotational equation gives
\begin{equation}
    \boldsymbol{\vartheta}_0
    =
    -
    \left(
        \mathbf{K}_{\vartheta\vartheta}
        -
        \omega^2I_{\mathrm{D}}\mathbf{I}
    \right)^{-1}
    \mathbf{K}_{\vartheta u}
    \mathbf{u}_0.
\end{equation}

For an acoustic mode,
\begin{equation}
    \omega\rightarrow0
    \qquad
    \text{as}
    \qquad
    k\rightarrow0,
\end{equation}
whereas the rotational gap remains finite,
\begin{equation}
    \omega_{\vartheta}^2(0)
    =
    \frac{
        16A_{\mathrm{ct}}\kappa_{\mathrm{c}}r_{\mathrm{D}}^2
    }{
        I_{\mathrm{D}}
    }.
\end{equation}

Hence
\begin{equation}
    \frac{\omega}{\omega_{\vartheta}(0)}
    \rightarrow0
\end{equation}
on the acoustic branch.

The exact frequency-dependent translational Schur complement is
\begin{equation}
    \mathbf{K}_{\mathrm{eff}}
    =
    \mathbf{K}_{uu}
    -
    \mathbf{K}_{u\vartheta}
    \left(
        \mathbf{K}_{\vartheta\vartheta}
        -
        \omega^2I_{\mathrm{D}}\mathbf{I}
    \right)^{-1}
    \mathbf{K}_{\vartheta u}.
\end{equation}

Using
\begin{equation}
    [\mathbf{k}]_{\times}^{\dagger}
    [\mathbf{k}]_{\times}
    =
    k^2\mathbf{I}
    -
    \mathbf{k}\mathbf{k}^{\mathsf{T}},
\end{equation}
one obtains the leading acoustic stiffness
\begin{equation}
    \boxed{
    \mathbf{K}_{\mathrm{eff}}^{(2)}
    =
    4A_{\mathrm{ct}}\kappa_{\mathrm{c}}r_{\mathrm{D}}^2
    \left(
        k^2\mathbf{I}
        +
        \mathbf{k}\mathbf{k}^{\mathsf{T}}
    \right).
    }
\end{equation}

For a transverse polarization,
\begin{equation}
    \mathbf{k}\cdot\mathbf{e}_{\mathrm{T}}=0,
\end{equation}
the stiffness eigenvalue is
\begin{equation}
    K_{\mathrm{T}}^{(2)}
    =
    4A_{\mathrm{ct}}\kappa_{\mathrm{c}}r_{\mathrm{D}}^2k^2.
\end{equation}

For a longitudinal polarization,
\begin{equation}
    \mathbf{e}_{\mathrm{L}}\parallel\mathbf{k},
\end{equation}
the stiffness eigenvalue is
\begin{equation}
    K_{\mathrm{L}}^{(2)}
    =
    8A_{\mathrm{ct}}\kappa_{\mathrm{c}}r_{\mathrm{D}}^2k^2.
\end{equation}

Thus the unconstrained isotropic scaffold satisfies
\begin{equation}
    c_{\mathrm{L}}^2
    =
    2c_{\mathrm{T}}^2.
\end{equation}

The dynamically relaxed microrotation may also be reconstructed explicitly. Since $[\mathbf{k}]_{\times}$ is real and antisymmetric,
\begin{equation}
    [\mathbf{k}]_{\times}^{\dagger}
    =
    -[\mathbf{k}]_{\times}.
\end{equation}

The leading Hermitian-conjugate block is therefore
\begin{equation}
    \mathbf{K}_{\vartheta u}^{(1)}
    =
    -8iA_{\mathrm{ct}}\kappa_{\mathrm{c}}r_{\mathrm{D}}^2[\mathbf{k}]_{\times}.
\end{equation}

Substitution into the rotational equation gives
\begin{equation}
    \boldsymbol{\vartheta}_0
    =
    \frac{i}{2}
    \mathbf{k}\times\mathbf{u}_0
    +
    \mathcal{O}\!\left(
        k^3d_{\mathrm{nn}}^2
        |\mathbf{u}_0|
    \right).
\end{equation}

With the Fourier convention
\begin{equation}
    \mathbf{u}
    \propto
    \exp
    \left[
        i
        \left(
            \mathbf{k}\cdot\mathbf{x}
            -
            \omega t
        \right)
    \right],
\end{equation}
one has
\begin{equation}
    \nabla\times\mathbf{u}
    =
    i\mathbf{k}\times\mathbf{u}.
\end{equation}

Therefore,
\begin{equation}
    \boxed{
    \boldsymbol{\vartheta}
    =
    \frac{1}{2}
    \nabla\times\mathbf{u}
    }
\end{equation}
at leading long-wavelength order.

The affine result $\boldsymbol{\vartheta}=\boldsymbol{\omega}$ is therefore recovered dynamically as the low-energy limit of the gapped microrotational sector rather than imposed as a static wave assumption.


\section{Isotropic Contact Specialization of the Active FCC Scaffold}
\label{sec:fcc_isotropy}

For a cubic elastic continuum, the anisotropy factor is
\begin{equation}
    \mathcal{A}_{\mathrm{cub}} = \frac{2C_{44}}{C_{11}-C_{12}}.
\end{equation}

Using the FCC constants derived above,
\begin{equation}
    \mathcal{A}_{\mathrm{cub}} = \frac{2\left(\kappa_{\mathrm{n}}+\kappa_{\mathrm{t}}\right)}{\kappa_{\mathrm{n}}+3\kappa_{\mathrm{t}}}.
\end{equation}

Exact isotropy requires $\mathcal{A}_{\mathrm{cub}}=1$, which gives
\begin{equation}
    \boxed{
    \kappa_{\mathrm{n}} = \kappa_{\mathrm{t}}.
    }
\end{equation}

Thus FCC geometry guarantees a cubic continuum, while equality of the incremental normal and tangential contact coefficients is the additional constitutive condition required for exact leading long-wavelength isotropy in the present nearest-neighbor parent model. This equality is not implied by FCC geometry alone. A deeper microscopic TFT completion must either derive it or replace the present contact structure by another interaction having the same isotropic continuum limit.

Define
\begin{equation}
    \kappa_{\mathrm{c}} \equiv \kappa_{\mathrm{n}} = \kappa_{\mathrm{t}},
\end{equation}
and define the baseline microscopic scaffold tangent shear coefficient by
\begin{equation}
    K_{\mathrm{sh}}^{(\mathrm{micro},0)} \equiv \frac{A_{\mathrm{ct}}d_{\mathrm{nn}}^2}{V_{\mathrm{p}}}\kappa_{\mathrm{c}}.
\end{equation}

Because the normal and tangential deformations are orthogonal projections of the same objective contact-point displacement,
\begin{equation}
    \left(\delta_{\mathrm{n}}^{(ij)}\right)^2
    +
    \left|\boldsymbol{\delta}_{\mathrm{t}}^{(ij)}\right|^2
    =
    \left|\Delta\mathbf{u}_{ij}^{\mathrm{c}}\right|^2.
\end{equation}

The translational contact energy therefore specializes to
\begin{equation}
    \mathcal{U}_{\mathrm{tr}}^{(2)}
    =
    \frac{1}{2}
    \kappa_{\mathrm{c}}
    \left|\Delta\mathbf{u}_{ij}^{\mathrm{c}}\right|^2.
\end{equation}

This compact one-coefficient form is a consequence of the isotropic specialization; it is not an independent microscopic derivation of $\kappa_{\mathrm{n}}=\kappa_{\mathrm{t}}$.

The isotropic scaffold constants are
\begin{equation}
    C_{11}^{(\mathrm{FCC})} = 2K_{\mathrm{sh}}^{(\mathrm{micro},0)}, \qquad C_{12}^{(\mathrm{FCC})} = 0, \qquad C_{44}^{(\mathrm{FCC})} = K_{\mathrm{sh}}^{(\mathrm{micro},0)}.
\end{equation}

The scaffold tensor is
\begin{equation}
    C_{ijkl}^{(\mathrm{FCC})} = K_{\mathrm{sh}}^{(\mathrm{micro},0)}\left(\delta_{ik}\delta_{jl}+\delta_{il}\delta_{jk}\right).
\end{equation}

Thus $K_{\mathrm{sh}}^{(\mathrm{micro},0)}$ is the unstrained primitive-cell tangent shear modulus of the FCC scaffold, homogenized over the full primitive-cell volume. It is a constitutive stiffness of the scaffold and is not itself a static or dynamic shear stress.

The scaffold has first Lam\'e coefficient zero and shear modulus $K_{\mathrm{sh}}^{(\mathrm{micro},0)}$. It still contributes the finite isotropic-compression response
\begin{equation}
    K_{\mathrm{bulk}}^{(\mathrm{FCC})} = \frac{2}{3}K_{\mathrm{sh}}^{(\mathrm{micro},0)}.
\end{equation}

Before the separate coherent interstitial-volume constraint is imposed, the dynamically reduced isotropic scaffold contains two degenerate transverse acoustic polarizations and one longitudinal acoustic polarization satisfying
\begin{equation}
    c_{\mathrm{L}}^2 = 2c_{\mathrm{T}}^2.
\end{equation}

The interstitial sector therefore does not generate the leading isotropy of the FCC scaffold. Its distinct coherent role is the microscopic cell-volume constraint developed in \hyperref[app:condensate_impedance]{Appendix~\ref*{app:condensate_impedance}}.


\section{Effective Modal Normalization of the FCC Contact Coefficient}

The preceding homogenization expresses the baseline microscopic scaffold tangent shear coefficient as
\begin{equation}
    K_{\mathrm{sh}}^{(\mathrm{micro},0)} = \frac{A_{\mathrm{ct}}d_{\mathrm{nn}}^2}{V_{\mathrm{p}}}\kappa_{\mathrm{c}}.
\end{equation}

The directional transverse contact law gives
\begin{equation}
    \kappa_{\mathrm{c}} = \frac{\tau_{\mathrm{D}}}{d_{\mathrm{nn}}^2}.
\end{equation}

Therefore,
\begin{equation}
    K_{\mathrm{sh}}^{(\mathrm{micro},0)} = \frac{A_{\mathrm{ct}}}{V_{\mathrm{p}}}\tau_{\mathrm{D}}.
\end{equation}

The quantity $A_{\mathrm{ct}}$ must be interpreted carefully. It is not introduced here as a directly observed geometric patch on a Wigner--Seitz face. It is an effective modal area that converts the surface-inertia and surface-tension representation of the scaffold contact action into the primitive-cell translational normalization.

The Wigner--Seitz cell of an FCC lattice is a rhombic dodecahedron. Its twelve faces correspond to the twelve nearest-neighbor directions. For nearest-neighbor separation $d_{\mathrm{nn}}$, each rhombic face has diagonals $d_{\mathrm{nn}}$ and $d_{\mathrm{nn}}/\sqrt{2}$. The full Wigner--Seitz face area is therefore
\begin{equation}
    A_{\mathrm{WS}} = \frac{1}{2}d_{\mathrm{nn}}\frac{d_{\mathrm{nn}}}{\sqrt{2}} = \frac{d_{\mathrm{nn}}^2}{2\sqrt{2}}.
\end{equation}

Using $d_{\mathrm{nn}}=2r_{\mathrm{D}}$, one obtains
\begin{equation}
    A_{\mathrm{WS}} = \sqrt{2}r_{\mathrm{D}}^2.
\end{equation}

The ratio of one full Wigner--Seitz face to the primitive-cell volume is
\begin{equation}
    \frac{A_{\mathrm{WS}}}{V_{\mathrm{p}}} = \frac{\sqrt{2}r_{\mathrm{D}}^2}{4\sqrt{2}r_{\mathrm{D}}^3} = \frac{1}{4r_{\mathrm{D}}}.
\end{equation}

This is a valid FCC geometric identity. It defines the full nearest-neighbor face area. It does not by itself identify the modal area associated with the reduced transverse scaffold transfer coordinate.

The modal area is fixed by matching the scaffold surface-inertia representation of the reduced transverse transfer coordinate to the primitive-cell nodal inertia normalization. Let $M_{\mathrm{D}}$ be the independently defined Soft Floor node inertia scale, and let $\mu_{\mathrm{D}}$ be the corresponding baseline scaffold surface inertia. The modal area is defined by
\begin{equation}
    A_{\mathrm{ct}}^{(\mathrm{modal})} \equiv \frac{M_{\mathrm{D}}}{\mu_{\mathrm{D}}}.
\end{equation}

This is a normalization identity, not a literal statement that each bond carries an independent mass $M_{\mathrm{D}}$. The primitive cell carries the shared nodal translational inertia once. For the reduced transverse scaffold coordinate used in the homogenized coefficient, the effective area $A_{\mathrm{ct}}$ is therefore identified with $A_{\mathrm{ct}}^{(\mathrm{modal})}$. This identification is constitutive and modal rather than a claim that the same area has been geometrically resolved as a physical contact patch on each FCC face.

Using the established TFT relations
\begin{equation}
    M_{\mathrm{D}} = \frac{2\pi\hbar}{c_{\mathrm{E}}r_{\mathrm{D}}},
\end{equation}
and
\begin{equation}
    \mu_{\mathrm{D}} = \frac{c_{\mathrm{E}}^2}{G r_{\mathrm{D}}},
\end{equation}
one obtains
\begin{equation}
    A_{\mathrm{ct}}^{(\mathrm{modal})} = \frac{2\pi\hbar/(c_{\mathrm{E}}r_{\mathrm{D}})}{c_{\mathrm{E}}^2/(G r_{\mathrm{D}})} = \frac{2\pi\hbar G}{c_{\mathrm{E}}^3}.
\end{equation}

Since the substrate scale satisfies
\begin{equation}
    \ell_{\mathrm{S}}^2 = \frac{\hbar G}{c_{\mathrm{E}}^3},
\end{equation}
the modal area is
\begin{equation}
    A_{\mathrm{ct}}^{(\mathrm{modal})} = 2\pi\ell_{\mathrm{S}}^2.
\end{equation}

The corresponding transverse scaffold modal projection is defined by
\begin{equation}
    \iota_{\mathrm{tr}}^{(\mathrm{modal})} \equiv \frac{A_{\mathrm{ct}}^{(\mathrm{modal})}}{V_{\mathrm{p}}}.
\end{equation}

Using
\begin{equation}
    V_{\mathrm{p}} = 4\sqrt{2}r_{\mathrm{D}}^3,
\end{equation}
and
\begin{equation}
    A_{\mathrm{ct}}^{(\mathrm{modal})} = 2\pi\ell_{\mathrm{S}}^2,
\end{equation}
one obtains
\begin{equation}
    \iota_{\mathrm{tr}}^{(\mathrm{modal})} = \frac{\pi\ell_{\mathrm{S}}^2}{2\sqrt{2}r_{\mathrm{D}}^3}.
\end{equation}

This projection also satisfies
\begin{equation}
    \iota_{\mathrm{tr}}^{(\mathrm{modal})} = \frac{\rho_{\mathrm{sh}}^{(\mathrm{micro},0)}}{\mu_{\mathrm{D}}},
\end{equation}
where
\begin{equation}
    \rho_{\mathrm{sh}}^{(\mathrm{micro},0)} \equiv \frac{M_{\mathrm{D}}}{V_{\mathrm{p}}}
\end{equation}
is the baseline primitive-cell FCC scaffold inertia density.

Indeed, since
\begin{equation}
    A_{\mathrm{ct}}^{(\mathrm{modal})} = \frac{M_{\mathrm{D}}}{\mu_{\mathrm{D}}},
\end{equation}
it follows that
\begin{equation}
    \frac{A_{\mathrm{ct}}^{(\mathrm{modal})}}{V_{\mathrm{p}}} = \frac{\rho_{\mathrm{sh}}^{(\mathrm{micro},0)}}{\mu_{\mathrm{D}}}.
\end{equation}

The quantity $\iota_{\mathrm{tr}}^{(\mathrm{modal})}$ is the transverse projection associated with the periodic FCC scaffold transfer mode. It is not the localized defect-flow projection $\iota_{\mathrm{V}}^{(\mathrm{def})}$ introduced in \hyperref[app:condensate_impedance]{Appendix~\ref*{app:condensate_impedance}}. The former belongs to the periodic scaffold contact normalization; the latter belongs to a boundary-value problem for relative condensate flow around a localized defect.

The equivalent circular radius associated with this modal area is
\begin{equation}
    r_{\mathrm{ct}}^{(\mathrm{modal})} = \sqrt{\frac{A_{\mathrm{ct}}^{(\mathrm{modal})}}{\pi}} = \sqrt{2}\ell_{\mathrm{S}}.
\end{equation}

This result should be read as follows: $A_{\mathrm{ct}}^{(\mathrm{modal})}$ is the surface-to-cell normalization area of the reduced transverse scaffold transfer coordinate. It should not be read as having been geometrically located as a literal physical patch on every FCC face.

The ratio to the full Wigner--Seitz face may still be useful as a dimensionless comparison:
\begin{equation}
    f_{\mathrm{ct}}^{(\mathrm{modal})} = \frac{A_{\mathrm{ct}}^{(\mathrm{modal})}}{A_{\mathrm{WS}}} = \sqrt{2}\pi \left( \frac{\ell_{\mathrm{S}}}{r_{\mathrm{D}}} \right)^2.
\end{equation}

Because $\ell_{\mathrm{S}} \ll r_{\mathrm{D}}$,
\begin{equation}
    f_{\mathrm{ct}}^{(\mathrm{modal})} \ll 1.
\end{equation}

This ratio is an effective modal-area ratio. It does not prove that only this fraction of the Wigner--Seitz face is physically active as a literal geometric contact footprint. The full Wigner--Seitz face remains the geometric nearest-neighbor boundary; the modal area is the normalization measure of the reduced transverse scaffold transfer coordinate.

The integrated modal stiffness associated with the reduced transfer coordinate is
\begin{equation}
    k_{\mathrm{ct}}^{(\mathrm{modal})} = A_{\mathrm{ct}}^{(\mathrm{modal})}\kappa_{\mathrm{c}}.
\end{equation}

Using
\begin{equation}
    \kappa_{\mathrm{c}} = \frac{\tau_{\mathrm{D}}}{d_{\mathrm{nn}}^2},
\end{equation}
one obtains
\begin{equation}
    k_{\mathrm{ct}}^{(\mathrm{modal})} = \frac{\pi}{2}\tau_{\mathrm{D}} \left( \frac{\ell_{\mathrm{S}}}{r_{\mathrm{D}}} \right)^2.
\end{equation}

The discrete directional speed associated with the reduced scaffold contact sequence is
\begin{equation}
    c_{\mathrm{T}}^{(0)} = d_{\mathrm{nn}} \sqrt{\frac{k_{\mathrm{ct}}^{(\mathrm{modal})}}{m_{\mathrm{ct}}^{(\mathrm{modal})}}}.
\end{equation}

The corresponding scaffold modal inertia is
\begin{equation}
    m_{\mathrm{ct}}^{(\mathrm{modal})} = A_{\mathrm{ct}}^{(\mathrm{modal})}\mu_{\mathrm{D}} = M_{\mathrm{D}}.
\end{equation}

This equality is a modal normalization of the representative scaffold transfer coordinate, not a statement that each of the six primitive-cell bonds contains a separate full node mass.

Using $k_{\mathrm{ct}}^{(\mathrm{modal})}$, $m_{\mathrm{ct}}^{(\mathrm{modal})}$, and $\kappa_{\mathrm{c}}$, one finds
\begin{equation}
\begin{aligned}
    \left(c_{\mathrm{T}}^{(0)}\right)^2 &= d_{\mathrm{nn}}^2 \frac{A_{\mathrm{ct}}^{(\mathrm{modal})}\kappa_{\mathrm{c}}}{A_{\mathrm{ct}}^{(\mathrm{modal})}\mu_{\mathrm{D}}} \\
    &= \frac{\tau_{\mathrm{D}}}{\mu_{\mathrm{D}}} = c_{\mathrm{E}}^2.
\end{aligned}
\end{equation}

Therefore,
\begin{equation}
    c_{\mathrm{T}}^{(0)} = c_{\mathrm{E}}.
\end{equation}

The same normalization follows from the dynamically reduced three-dimensional transverse branch. Using
\begin{equation}
    K_{\mathrm{T}}^{(2)}
    =
    4A_{\mathrm{ct}}\kappa_{\mathrm{c}}r_{\mathrm{D}}^2k^2
\end{equation}
and
\begin{equation}
    d_{\mathrm{nn}}^2=4r_{\mathrm{D}}^2,
\end{equation}
the transverse speed is
\begin{equation}
\begin{aligned}
    c_{\mathrm{T}}^2
    &=
    \frac{
        4A_{\mathrm{ct}}\kappa_{\mathrm{c}}r_{\mathrm{D}}^2
    }{
        M_{\mathrm{D}}
    }\\
    &=
    \frac{
        A_{\mathrm{ct}}d_{\mathrm{nn}}^2\kappa_{\mathrm{c}}
    }{
        M_{\mathrm{D}}
    }\\
    &=
    \frac{\tau_{\mathrm{D}}}{\mu_{\mathrm{D}}}
    =
    c_{\mathrm{E}}^2.
\end{aligned}
\end{equation}

Thus the directional normalization and the dynamically relaxed three-dimensional FCC transverse branch yield the same unstrained characteristic speed.

The modal area cancels from the unstrained wave-speed normalization. Its role is not to fit $c_{\mathrm{E}}$, but to keep the surface and primitive-cell representations of the scaffold transfer coordinate mutually consistent.

The baseline primitive-cell scaffold inertia density follows directly from the node mass divided by the full primitive-cell volume:
\begin{equation}
    \rho_{\mathrm{sh}}^{(\mathrm{micro},0)} \equiv \frac{M_{\mathrm{D}}}{V_{\mathrm{p}}}.
\end{equation}

Using
\begin{equation}
    M_{\mathrm{D}} = \rho_{\mathrm{D}} \left( \frac{4}{3}\pi r_{\mathrm{D}}^3 \right),
\end{equation}
and
\begin{equation}
    V_{\mathrm{p}} = 4\sqrt{2}r_{\mathrm{D}}^3,
\end{equation}
one obtains
\begin{equation}
    \rho_{\mathrm{sh}}^{(\mathrm{micro},0)} = \frac{\pi}{3\sqrt{2}}\rho_{\mathrm{D}}.
\end{equation}

The FCC close-packing fraction is
\begin{equation}
    \beta_{\mathrm{geo}} = \frac{\pi}{3\sqrt{2}}.
\end{equation}

Thus,
\begin{equation}
    \rho_{\mathrm{sh}}^{(\mathrm{micro},0)} = \beta_{\mathrm{geo}}\rho_{\mathrm{D}}.
\end{equation}

Here $\rho_{\mathrm{D}}$ is the density assigned to the solid nodal phase itself, whereas $\rho_{\mathrm{sh}}^{(\mathrm{micro},0)}$ is the corresponding nodal mass averaged over the complete FCC primitive-cell volume. The factor $\beta_{\mathrm{geo}}$ therefore enters because only the fraction $\beta_{\mathrm{geo}}$ of the primitive-cell volume is occupied by the solid FCC nodal phase. It is a geometric volume-averaging factor and does not represent a reduction of the intrinsic nodal density $\rho_{\mathrm{D}}$.

The corresponding baseline primitive-cell scaffold tangent shear coefficient is
\begin{equation}
    K_{\mathrm{sh}}^{(\mathrm{micro},0)} = \rho_{\mathrm{sh}}^{(\mathrm{micro},0)}c_{\mathrm{E}}^2 = \beta_{\mathrm{geo}}\rho_{\mathrm{D}}c_{\mathrm{E}}^2.
\end{equation}

Equivalently,
\begin{equation}
    K_{\mathrm{sh}}^{(\mathrm{micro},0)} = \frac{A_{\mathrm{ct}}^{(\mathrm{modal})}}{V_{\mathrm{p}}}\tau_{\mathrm{D}}.
\end{equation}

The coefficient $K_{\mathrm{sh}}^{(\mathrm{micro},0)}$ is the unstrained primitive-cell tangent shear modulus of the FCC scaffold, homogenized over the full primitive-cell volume. It is a constitutive stiffness and is not itself a static shear stress, a gravitational-wave stress, a measure of local gravitational loading, a localized defect impedance coefficient, or an effective macroscopic gravitational-wave response coefficient.

The associated inverse-length projection is $\iota_{\mathrm{tr}}^{(\mathrm{modal})}$, not the localized defect projection $\iota_{\mathrm{V}}^{(\mathrm{def})}$. The scaffold coefficient must therefore not be identified either with the defect-flow impedance calibration or with an effective macroscopic gravitational-wave response coefficient.

When the selected coherent acoustic coordinate is discussed in \hyperref[app:condensate_impedance]{Appendix~\ref*{app:condensate_impedance}}, the corresponding unstrained kinetic coefficient may be denoted $\rho_{\mathrm{coh}}^{(0)}$ to emphasize the retained coherent coordinate. At the primitive-cell level established here,
\begin{equation}
    \rho_{\mathrm{coh}}^{(0)}
    \equiv
    \rho_{\mathrm{sh}}^{(\mathrm{micro},0)}
    =
    \frac{M_{\mathrm{D}}}{V_{\mathrm{p}}}.
\end{equation}

This notation does not establish an independently additive interstitial common-mode inertia. A deeper coupled scaffold--condensate kinetic completion, if later derived, must modify the coherent kinetic coefficient consistently rather than be added by double counting the scaffold normalization established here.


% BEGIN PAPER-IV SYNCHRONIZATION: constitutive diagnostic
\subsection*{Nonlinear phase--deformation coupling is not fixed by the tangent action}
\phantomsection\label{sec:D_phase_deformation_diagnostic}

The diagonal reference Hessian fixes small-amplitude contact response, not every nonlinear coupling of a localized state. This distinction can be demonstrated without adopting another TFT interaction. For one bond, define the dimensionless normal deformation $s=\delta_{\mathrm n}/r_{\mathrm D}$ and consider the diagnostic family of integrated energies
\[
 U_g=A_{\mathrm{ct}}r_{\mathrm D}^{2}
 \left\{\frac{\kappa_{\mathrm n}}2
 \left[s-g(1-\cos\Delta\varphi)\right]^2
 +\kappa_\varphi(1-\cos\Delta\varphi)\right\}.
\]
Here $g$ is an arbitrary dimensionless comparison parameter, not a proposed or calibrated TFT constant. For positive $\kappa_{\mathrm n},\kappa_\varphi$, every member is nonnegative, periodic in $\Delta\varphi$, and even under reversal of that phase difference. The reference quadratic energy and Hessian are independent of $g$:
\[
 U_g^{(2)}=\frac{A_{\mathrm{ct}}r_{\mathrm D}^{2}}2
 (\kappa_{\mathrm n}s^2+\kappa_\varphi\Delta\varphi^2),
 \qquad
 \left.\nabla^2_{(s,\Delta\varphi)}U_g\right|_0
 =A_{\mathrm{ct}}r_{\mathrm D}^{2}
 \begin{pmatrix}\kappa_{\mathrm n}&0\\0&\kappa_\varphi\end{pmatrix}.
\]
Nevertheless, minimizing over $s$ at fixed phase gives
\[
 s_*=g(1-\cos\Delta\varphi),\qquad
 U_g(s_*,\Delta\varphi)
 =A_{\mathrm{ct}}\kappa_\varphi r_{\mathrm D}^{2}
 (1-\cos\Delta\varphi),
\]
while
\[
 \left.\frac{\partial^3 U_g}
 {\partial s\,\partial(\Delta\varphi)^2}\right|_0
 =-A_{\mathrm{ct}}r_{\mathrm D}^{2}\kappa_{\mathrm n}g.
\]
Thus even identical reference tangent coefficients and identical relaxed phase potentials need not determine identical phase-induced boundary motion. The diagnostic compares local bond energies; independently minimizing all bonds is not a construction of a geometrically compatible lattice defect.

Likewise, eliminating a stable linear shape coordinate from
$U=U_\varphi+\tfrac12 k_q q^2-qF_\varphi$, with $k_q>0$,
gives $U_{\mathrm{relaxed}}=U_\varphi-F_\varphi^2/(2k_q)$.
The relaxed term lowers the energy. It does not, by itself, supply an added positive localization penalty. A stable finite-amplitude phase--surface configuration must be tested using its complete nonlinear energy, admissible geometry, and kinetic metric. These diagnostics leave the reference action and its spectrum unchanged.
% END PAPER-IV SYNCHRONIZATION

\section{Status of the FCC Dynamic Spectrum and Remaining Closures}

This appendix establishes the baseline contact mechanics, explicit harmonic Bloch operator, and long-wavelength acoustic reduction of the active FCC Soft Floor. The results define the gravity-independent FCC scaffold carrier but do not constitute a complete ultraviolet theory of every physical branch.

The first completed result is the baseline nearest-neighbor directional transverse scaffold action,
\begin{equation}
    \omega^2(k) = \frac{4c_{\mathrm{E}}^2}{d_{\mathrm{nn}}^2} \sin^2 \left( \frac{kd_{\mathrm{nn}}}{2} \right),
\end{equation}
with long-wavelength limit
\begin{equation}
    c_{\mathrm{T}}^{(0)} = c_{\mathrm{E}}.
\end{equation}

This scalar chain establishes the nearest-neighbor transverse normalization only. It is not the complete three-dimensional FCC photon dispersion. The selected positive dispersion-suppressed photon branch is developed separately in \hyperref[app:photon_dispersion_relaxation]{Appendix~\ref*{app:photon_dispersion_relaxation}}. Its effective higher-gradient coefficient $1/12$ is not derived here from the nearest-neighbor parent action.

The second completed result is the objective nearest-neighbor contact kinematics. The local contact deformation is determined by the contact-point displacement
\begin{equation}
    \Delta\mathbf{u}_{ij}^{\mathrm{c}} = \mathbf{u}_j - \mathbf{u}_i - r_{\mathrm{D}} \left( \boldsymbol{\vartheta}_i + \boldsymbol{\vartheta}_j \right) \times \hat{\mathbf{n}}_{ij},
\end{equation}
rather than by the relative center displacement alone. This construction yields mechanical contact measures that vanish under rigid translation and compatible rigid rotation. The extended internal phase difference is invariant under a uniform phase shift.

The third completed result is the parity-even quadratic contact action. The baseline active vacuum admits normal, tangential, bending, twisting, and phase-difference quadratic terms. It excludes the tangential--bending pseudoscalar unless an explicit chiral pseudoscalar is introduced. The unpolarized active vacuum therefore contains no baseline term proportional to
\begin{equation}
    \boldsymbol{\delta}_{\mathrm{t}}^{(ij)} \cdot \boldsymbol{\delta}_{\mathrm{b}}^{(ij)}.
\end{equation}

The fourth completed result is the primitive-cell bookkeeping. The six undirected nearest-neighbor bonds contribute potential energy once, while the translational mass and microrotational inertia of the retained acoustic parent are assigned once through
\begin{equation}
    \mathbf{M}_{\mathrm{cell}}^{(6)}
    =
    \operatorname{diag}
    \left(
        M_{\mathrm{D}}\mathbf{I}_3,
        I_{\mathrm{D}}\mathbf{I}_3
    \right).
\end{equation}
No full seven-component dynamical determinant is asserted because an independent pristine-vacuum phase inertia has not been derived.

The fifth completed result is the explicit FCC Bloch stiffness,
\begin{equation}
    \mathbf{K}_{\mathrm{FCC}}^{(6)}(\mathbf{k})
    =
    A_{\mathrm{ct}}
    \sum_{a=1}^{6}
    \mathbf{B}_a^{(6)\dagger}
    \mathbf{H}_{\mathrm{c}}^{(6)}
    \mathbf{B}_a^{(6)}.
\end{equation}

For positive retained tangent coefficients,
\begin{equation}
    \mathbf{K}_{\mathrm{FCC}}^{(6)}(\mathbf{k})\succeq0
\end{equation}
for every $\mathbf{k}$. With positive primitive-cell inertia, this implies
\begin{equation}
    \omega_s^2(\mathbf{k})\geq0
\end{equation}
for every branch $s$ of the retained six-degree-of-freedom harmonic parent system.

The sixth completed result is the finite microrotational zone-center gap,
\begin{equation}
    \omega_{\vartheta}^2(0)
    =
    \frac{
        16A_{\mathrm{ct}}\kappa_{\mathrm{t}}r_{\mathrm{D}}^2
    }{
        I_{\mathrm{D}}
    }.
\end{equation}
The rotational sector is therefore optical rather than acoustic.

The seventh completed result is the dynamic low-energy Schur reduction. Under the isotropic contact specialization,
\begin{equation}
    \kappa_{\mathrm{n}}=\kappa_{\mathrm{t}}=\kappa_{\mathrm{c}},
\end{equation}
the leading acoustic translational stiffness becomes
\begin{equation}
    \mathbf{K}_{\mathrm{eff}}^{(2)}
    =
    4A_{\mathrm{ct}}\kappa_{\mathrm{c}}r_{\mathrm{D}}^2
    \left(
        k^2\mathbf{I}
        +
        \mathbf{k}\mathbf{k}^{\mathsf{T}}
    \right),
\end{equation}
and the relaxed microrotation obeys
\begin{equation}
    \boldsymbol{\vartheta}
    =
    \frac12\nabla\times\mathbf{u}
\end{equation}
to leading long-wavelength order.

The eighth completed result is the affine FCC homogenization. The six nearest-neighbor directions produce a cubic continuum with
\begin{equation}
    C_{11}^{(\mathrm{FCC})} = \frac{A_{\mathrm{ct}}d_{\mathrm{nn}}^2}{V_{\mathrm{p}}} \left( \kappa_{\mathrm{n}} + \kappa_{\mathrm{t}} \right),
\end{equation}
\begin{equation}
    C_{12}^{(\mathrm{FCC})} = \frac{A_{\mathrm{ct}}d_{\mathrm{nn}}^2}{2V_{\mathrm{p}}} \left( \kappa_{\mathrm{n}} - \kappa_{\mathrm{t}} \right),
\end{equation}
and
\begin{equation}
    C_{44}^{(\mathrm{FCC})} = \frac{A_{\mathrm{ct}}d_{\mathrm{nn}}^2}{2V_{\mathrm{p}}} \left( \kappa_{\mathrm{n}} + \kappa_{\mathrm{t}} \right).
\end{equation}

FCC geometry guarantees cubic symmetry. Exact leading isotropy follows only under the constitutive specialization
\begin{equation}
    \kappa_{\mathrm{n}} = \kappa_{\mathrm{t}}.
\end{equation}

The ninth completed result is the effective scaffold modal normalization,
\begin{equation}
    A_{\mathrm{ct}}^{(\mathrm{modal})} = \frac{M_{\mathrm{D}}}{\mu_{\mathrm{D}}} = 2\pi\ell_{\mathrm{S}}^2.
\end{equation}

This quantity is a surface-to-cell normalization area for the reduced transverse scaffold transfer coordinate. It is not interpreted as a directly established literal contact patch on each Wigner--Seitz face.

The tenth completed result is the baseline primitive-cell scaffold density and tangent shear normalization,
\begin{equation}
    \rho_{\mathrm{sh}}^{(\mathrm{micro},0)} = \beta_{\mathrm{geo}}\rho_{\mathrm{D}},
    \qquad
    K_{\mathrm{sh}}^{(\mathrm{micro},0)} = \beta_{\mathrm{geo}}\rho_{\mathrm{D}}c_{\mathrm{E}}^2.
\end{equation}

Here $\rho_{\mathrm{D}}$ is the density of the solid nodal phase, whereas $\rho_{\mathrm{sh}}^{(\mathrm{micro},0)}$ is the corresponding nodal mass averaged over the complete FCC primitive-cell volume. The factor $\beta_{\mathrm{geo}}$ is therefore the geometric volume-averaging factor of the solid scaffold and does not represent a reduction of the intrinsic nodal density $\rho_{\mathrm{D}}$.

The coefficient $K_{\mathrm{sh}}^{(\mathrm{micro},0)}$ is the unstrained primitive-cell tangent shear modulus of that scaffold. It is a constitutive stiffness and is not itself a static shear stress, a gravitational-wave stress, a measure of local gravitational loading, a localized defect impedance coefficient, or an effective macroscopic gravitational-wave response coefficient.

Several microscopic completions remain outside the scope of this appendix.

The independent pristine-vacuum phase kinetic sector remains unresolved. Since the contact energy depends on $\Delta\varphi_{ij}$, a uniform phase displacement is a zero-energy direction of the phase potential. Whether the phase coordinate is gauge-redundant, locally locked, gapped, defect-confined, or dynamically active through an independently derived phase inertia is not determined by the present acoustic reduction. If a finite positive phase inertia is derived, the resulting additional branches must be analyzed explicitly.

Although harmonic non-negativity of the retained translation--microrotation parent is now established analytically, the complete classification of every finite-wave-number branch, degeneracy, avoided crossing, and polarization throughout the first Brillouin zone remains open.

The isotropic condition
\begin{equation}
    \kappa_{\mathrm{n}}=\kappa_{\mathrm{t}}
\end{equation}
remains a constitutive specialization. Its microscopic origin must be derived by a deeper active-vacuum contact theory or replaced by a more complete interaction structure having the same isotropic continuum limit.

The coherent interstitial-volume constraint is not produced by the scaffold mechanics alone. Its discrete Wigner--Seitz-cell formulation, pressure reaction, acoustic Bloch constraint, and continuum reduction are developed separately in \hyperref[app:condensate_impedance]{Appendix~\ref*{app:condensate_impedance}}.

The positive centered-second-difference contribution with coefficient $1/12$ used in the reduced photon action of \hyperref[app:photon_dispersion_relaxation]{Appendix~\ref*{app:photon_dispersion_relaxation}} is not generated by the nearest-neighbor parent action derived here. The complete ultraviolet photon eigenbranch therefore remains an open microscopic problem.

The chiral extension is similarly deferred. Parity-odd terms may enter only with an explicit pseudoscalar chirality and must vanish in the unpolarized vacuum average.

The mapping from the primitive-cell FCC dynamics derived here to a macroscopic gravitational-wave tensor response is not part of the present homogenization. That problem requires a separate collective-mode projection and macroscopic response construction and is developed in \hyperref[app:kinematic_shear]{Appendix~\ref*{app:kinematic_shear}}.

The localized defect-induced relative-flow impedance is developed separately in \hyperref[app:condensate_impedance]{Appendix~\ref*{app:condensate_impedance}}. The gravitational endpoint-transfer map and its referred-inertia specialization are developed separately in \hyperref[app:gravitational_endpoint_transfer]{Appendix~\ref*{app:gravitational_endpoint_transfer}}.

The status of the present appendix is therefore precise: it establishes the baseline FCC scaffold contact action, objective bond kinematics, primitive-cell kinetic bookkeeping, explicit six-degree-of-freedom harmonic Bloch operator, harmonic spectral non-negativity, finite microrotational gap, dynamic acoustic Schur reduction, affine cubic acoustic tensor, isotropic contact specialization, modal normalization, and microscopic scaffold coefficient scale. It does not include the coherent interstitial-volume constraint itself, the independent phase-sector dynamics, the complete ultraviolet photon spectrum, the gravity-specific endpoint closure, the localized relative-condensate-flow solution, or the collective projection from microscopic FCC dynamics to the macroscopic gravitational-wave tensor sector.
